
\section*{Abstract}
\addcontentsline{toc}{section}{Abstract}
\vspace{1cm}

\large
Associating a combinatorial meaning to abstract numbers has
been a driving force in mathematics,
with recent works by several researchers demonstrating how counting complexity can unlock deep structural insights
\cite{pak_WhatCombinatorialInterpretation_2022,ikenmeyer_WhatWhatNot_2022}.
This paper builds upon their notion in order to study the combinatorial structure of the $\textbf{\#PPAD}-1$ complexity class
using the \textit{PureCircuit},
a circuit-like problem that uses three-valued logic to find satisfying assignments.
The main contribution builds upon the relation of \textit{PureCircuit} with topological structures
and its potential usages for parsimonious reductions across \textbf{PPAD} problems.

We first demonstrate how to embed the three-valued logic into the \textbf{PPAD} framework and
the different parsimonious reductions we can construct across the \textit{EndOfLine}, 
$\textit{SourceOrExcess}(2,1)$ and \textit{PureCircuit} problems.
These reductions allow us to analyse unbalanced digraphs through a topological lens,
uncovering previously unexplored combinatorial classes within $\textbf{\#PPAD}-1$.
Moreover, our analysis extends to multi-source graph configurations, where we characterise their combinatorial structures
and demonstrate their connections between combinatorial and topological interpretations.

We complement our theoretical work with computational tools for visualising \textit{PureCircuit} instances.
The paper demonstrates the application of several meta-heuristic algorithms to approximate solutions
as well as a specialised backtracking algorithm for solution enumeration,
demonstrating practical improvements over brute-force approaches.




\vspace{1cm}

\noindent \textit{Keywords: Counting complexity, Kleene logic, Combinatorial Interpretations,
    Search Problems, PPAD, Hazard-free Circuits, PureCircuit, Complexity Theory}